\documentclass[12pt]{article}
\usepackage[utf8]{inputenc}
\usepackage[spanish]{babel}
\usepackage{amsmath}
\usepackage{amsfonts}
\usepackage{amssymb}
\usepackage{geometry}
\geometry{letterpaper, margin=2.5cm}

\begin{document}

\begin{center}
    \Large \textbf{Evaluación de Examen 1: Unidad I} \\
    \large Física para Ciencias de la Salud (005-1713) \\
    \vspace{0.5cm}
    \textbf{Estudiante:} Daniela Vallenilla \quad \textbf{C.I.:} 32.778.547
    \vspace{0.3cm} \\
    \large \textbf{Calificación Final: 9/10 puntos}
\end{center}

\vspace{0.5cm}

\section*{Análisis de Resultados}

\subsection*{1. Ángulo plano (Conversión a radianes)}
\textbf{Procedimiento y resultado:} Aplica correctamente el factor de conversión multiplicando por $\frac{\pi}{180}$. Realiza una simplificación de fracciones paso a paso $\left(\frac{75}{180} = \frac{15}{36} = \frac{5}{12}\right)$ muy limpia y metódica, obteniendo el valor exacto de $\frac{5\pi}{12}$ rad. \\
\textbf{Veredicto:} \textbf{Correcto} (2/2 pts).

\subsection*{2. Sistemas de coordenadas (Longitud del hueso)}
\textbf{Procedimiento y resultado:} Utiliza la fórmula formal de la distancia entre dos puntos y sustituye las coordenadas de manera correcta: $L = \sqrt{(7-1)^2 + (10-2)^2}$. Resuelve las operaciones algebraicas impecablemente hasta llegar al resultado de $10$ cm. \\
\textbf{Veredicto:} \textbf{Correcto} (2/2 pts).

\subsection*{3. Vectores (Fuerza resultante)}
\textbf{Procedimiento y resultado:} La estudiante es muy ordenada al desglosar el cálculo: suma primero las componentes en el eje x ($\hat{i}$) y luego en el eje y ($\hat{j}$). Esto evita errores de signos y le permite ensamblar correctamente el vector de fuerza resultante final: $\vec{F}_R = 30\hat{i} + 60\hat{j}$ N. \\
\textbf{Veredicto:} \textbf{Correcto} (2/2 pts).

\subsection*{4. Potencias de diez (Notación científica y prefijo)}
\textbf{Procedimiento y resultado:} Transforma la medida de forma adecuada empleando potencias de base diez ($12,5 \times 10^{-6}$ m). Sin embargo, omite la respuesta final específica requerida en el enunciado, la cual consistía en identificar y escribir el nombre del prefijo del Sistema Internacional (micrómetros). \\
\textbf{Veredicto:} \textbf{Incompleto / Parcialmente correcto} (1/2 pts).

\subsection*{5. Sistemas de unidades (Conversión de densidad)}
\textbf{Procedimiento y resultado:} Plantea los factores de conversión en cadena de manera perfecta, empleando las relaciones de masa y volumen adecuadas. Simplifica el cálculo correctamente realizando la operación $1,03 \times 1000$, obteniendo el valor numérico exacto de $1030$. \\
\textbf{Veredicto:} \textbf{Correcto} (2/2 pts).

\vspace{0.5cm}
\section*{Comentarios Generales y Sugerencias}
¡Excelente examen! La estudiante demuestra un manejo sumamente claro del álgebra, orden en sus resoluciones y una simplificación de fracciones y ecuaciones muy destacable. 
\begin{itemize}
    \item Como pequeña sugerencia formativa, se recomienda tener mayor cuidado con la notación al escribir las unidades derivadas finales (como expresar claramente $\text{kg/m}^3$ en forma de cociente) y recordar leer bien para responder todas las partes del enunciado, como ocurrió con el prefijo en la Pregunta 4.
\end{itemize}

\end{document}